\documentclass[10pt,a4paper]{amsart}
\usepackage[margin=19mm]{geometry}
\usepackage{amsmath,amssymb,mathtools}
\usepackage[scr=rsfs,cal=euler]{mathalfa}
\usepackage{xfrac}
\usepackage[unicode,hidelinks,bookmarks=false]{hyperref}
\newcommand{\ie}{i.e.\ }
\newcommand{\cf}{cf.\ }
\newcommand{\resp}{respectively\ }
\newcommand{\Subsection}{\subsection}
\providecommand{\nobreakdash}{\nobreak\hbox{-}}
\setlength{\emergencystretch}{2em}
\multlinegap0pt
\usepackage{amssymb}
\usepackage{mathtools}
\usepackage{enumerate}
\usepackage{multirow}
\newtheorem{lemma}{Lemma}
\newcommand {\mfp}{{\mathfrak{P}}}
\title{Solving Fermat-type equations over quadratic fields}
\author{Begum Gulsah Cakti}
\date{Source: arXiv:2509.14280v2 (CC BY 4.0)}
\begin{document}
\maketitle

\noindent\emph{Source and licence.} Solving Fermat-type equations over quadratic fields. Version arXiv:2509.14280v2 (CC BY 4.0), distributed by arXiv under the Creative Commons Attribution 4.0 International licence, http://creativecommons.org/licenses/by/4.0/, as recorded in the arXiv licence metadata for this exact version and shown on the abstract page https://arxiv.org/abs/2509.14280v2. Full licence notice: \texttt{licenses/arxiv-2509.14280-CC-BY-4.0.txt}.\par\smallskip

\noindent\textit{Editorial scope.} Counted unit: the article's lemma Epotgoodred (source0.tex lines 402--404). The article's complete original proof of it is delivered with it (source0.tex lines 406--419, one \texttt{proof} environment of 14 lines). No mathematics is written, repaired or re-derived: the statement and the proof are the article's own lines, in the article's own notation, and no retained line is substituted or re-flowed.  No further source-local context is needed: every cross-reference of every retained line resolves inside the two delivered spans.  Nothing was omitted from any retained span, so the package carries no omission record.  No retained line contains a citation, so the wrapper carries no bibliography and no bibliography entry.  No retained line contains a figure environment, an image-inclusion command or drawing code, so no image is delivered and the package declares no asset dependency.  Editorial formatting only, in addition: an AMS \texttt{amsart} preamble that restates the article's own declarations and macros used by the retained lines, namely the environments \texttt{lemma} on separate counters, together with the standard packages those lines need.  The retained environments print in this document as Lemma 1 in order of appearance, whereas the article prints the same items with the numbering of its own full document.  The complete native source of the article as distributed in the arXiv e-print for this exact version is bundled unchanged as \texttt{sources/185229/source0.tex}.
\begin{lemma}\label{lemma: Epotgoodred}
    Let $d=-3,-11,-19,-43$ and assume that $2\nmid abc$. The elliptic curve $E$ has potentially good reduction at $\mfp$.
\end{lemma}

\begin{proof}
   Let $d=-3,-11,-19,-43$ and assume that $2\nmid abc$. Note that 
  \begin{eqnarray*}
    v_{\mfp}(c_4)&=&v_{\mfp}(2^4(b^{2p}-d^ra^pc^p))\\
    &=&4v_{\mfp}(2)+v_{\mfp}(b^{2p}-d^ra^pc^p)\\
    &=&4+v_\mfp(b^{2p}-d^ra^pc^p)\\
    &\ge& 5
\end{eqnarray*}and that
    \begin{eqnarray*}
    v_{\mfp}(\Delta_E)&=& v_{\mfp}(2^4d^{2r}(abc)^{2p})\\
    &=&4v_{\mfp}(2)+2rv_{\mfp}(d)+2p(v_{\mfp}(a)+v_{\mfp}(b)+v_{\mfp}(c))\\
    &=&4
    \end{eqnarray*}Therefore,$$v_{\mfp}(j_E)\ge 11>0$$and $E$ has potentially good reduction at $\mfp$.
\end{proof}

\end{document}
